\begin{tikzpicture}[figure picture,x=10mm,y=10mm,font=\figurefont\fontsize{8}{10}\selectfont]
\path[use as bounding box] (0,0) rectangle (16.4,6.6);
% 列标题
\node[figure heading,anchor=west] at (0.4,6.35) {(上) 确定映射：每个输入点整颗送到唯一输出点（推前测度）};
\node at (2.2,5.85) {输入空间 $X$};
\node at (9.6,5.85) {输出空间 $Y$};
% 上排：确定映射
\fill[FigureInk] (2.2,5.0) circle[radius=1mm]; \node[left] at (2.1,5.0) {$x_1$};
\fill[FigureInk] (2.2,4.25) circle[radius=1mm]; \node[left] at (2.1,4.25) {$x_2$};
\fill[FigureInk] (2.2,3.5) circle[radius=1mm]; \node[left] at (2.1,3.5) {$x_3$};
\fill[FigureInk] (9.6,5.0) circle[radius=1mm]; \node[right] at (9.7,5.0) {$T(x_1)$};
\fill[FigureInk] (9.6,3.875) circle[radius=1mm]; \node[right] at (9.7,3.875) {$T(x_2)=T(x_3)$};
\draw[FigureBlue,-{Stealth[length=1.8mm]},line width=.9pt] (2.35,4.98)--(9.45,5.0);
\draw[FigureBlue,-{Stealth[length=1.8mm]},line width=.9pt] (2.35,4.27)--(9.45,3.93);
\draw[FigureBlue,-{Stealth[length=1.8mm]},line width=.9pt] (2.35,3.48)--(9.45,3.82);
\node[figure note] at (6.0,3.05) {多个源点可合并到同一目标，但每个源点只去一个点};
% 下排：测度核
\node[figure heading,anchor=west] at (0.4,2.45) {(下) 测度核：一个输入点按权重把质量分散给多个输出点};
\node at (2.2,1.95) {输入空间 $X$};
\node at (9.6,1.95) {输出空间 $Y$};
\fill[FigureInk] (2.2,1.05) circle[radius=1mm]; \node[left] at (2.1,1.05) {$x$};
\fill[FigureInk] (9.6,1.45) circle[radius=1mm]; \node[right] at (9.7,1.45) {$y_1$};
\fill[FigureInk] (9.6,.55) circle[radius=1mm]; \node[right] at (9.7,.55) {$y_2$};
\draw[FigureRed,-{Stealth[length=1.8mm]},line width=.9pt] (2.35,1.09)--(9.45,1.42);
\draw[FigureRed,-{Stealth[length=1.8mm]},line width=.9pt] (2.35,1.01)--(9.45,.58);
\node[above,FigureRed] at (5.9,1.42) {$K(x,\{y_1\})$};
\node[below,FigureRed] at (5.9,.62) {$K(x,\{y_2\})$};
\end{tikzpicture}
